\section{Algorithms}
\label{sec:algorithms}

Three algorithms drive the proposed module's higher-level behavior.
The first is a tiled multi-QR decoder that recovers more than one QR
sticker from a single camera frame
(Section~\ref{sec:algo-qr}). The second is a transparent pixel-level
RGB classifier that estimates banana spotting percentage from a JPEG
capture (Section~\ref{sec:algo-banana}). The third is the per-product
risk-score model that combines all sensing channels into a single
0--100 freshness indicator (Section~\ref{sec:algo-risk}).

\subsection{Multi-QR Tiled Decoder}
\label{sec:algo-qr}

A naïve single-pass QR decoder applied to a single JPEG capture from
the ESP32-CAM typically recovers at most one QR symbol per frame,
because most open-source decoders return on the first successful
detection. In a smart-cabinet setting this is suboptimal: several
products can be visible simultaneously, and forcing the user to scan
them one at a time defeats the purpose of automatic registration.

We implement a multi-QR pipeline that operates on a single capture.
The pipeline tiles the frame into a $3 \times 3$ grid of overlapping
windows, decodes each tile independently, and accumulates the union of
the recovered payloads. Concretely:

\begin{enumerate}
  \item The captured JPEG is downscaled to 480~px on its longer edge
        to bound the decoder's per-pass runtime.
  \item Nine overlapping windows are produced: each window covers
        $\sim\!45\%$ of the frame width and height with a 25\% overlap
        between neighbors, which guarantees that any QR fitting in a
        $\sim\!45\%$ region is fully contained in at least one window.
  \item Each window is passed to a portable QR decoder (\texttt{zxing2}
        on the web target, \texttt{qr\_code\_scanner} on the Android
        target). The decoder returns zero or more payloads per window.
  \item Decoded payloads are deduplicated by parsed JSON
        \texttt{productId} and inserted into a session set, so that a
        sticker physically present in two overlapping windows is
        registered exactly once.
\end{enumerate}

The pipeline runs in less than 200~ms per frame on a midrange Android
device. We trade a small amount of redundant decoding work for
substantially higher coverage of multi-product scenes.

\subsection{Pixel-Level RGB Banana Ripeness Classifier}
\label{sec:algo-banana}

Visible browning is a strong qualitative indicator of late-stage
banana ripeness. Rather than train a CNN to classify ripeness ---
which would require a curated dataset, hardware-accelerated inference,
or a cloud round-trip --- we implement a transparent pixel-level RGB
analysis that runs in a few milliseconds on commodity hardware and
returns a single interpretable percentage.

Given a JPEG capture $I$ in sRGB space, the classifier first segments
the banana surface from the background:
\begin{align*}
\mathcal{B}_{\text{yellow}} &= \{(r, g, b) :
    r > 150 \land g > 130 \land b < 110\} \\
\mathcal{B}_{\text{brown}}  &= \{(r, g, b) :
    100 \leq r \leq 200 \land 50 \leq g \leq 130
    \land b < 90\} \\
\mathcal{B}_{\text{dark}}   &= \{(r, g, b) :
    r < 90 \land g < 80 \land b < 70\} \\
\mathcal{B}_{\text{bg}}     &= \{(r, g, b) : r > 220 \land g > 220
    \land b > 220\}.
\end{align*}

Each pixel of $I$ is assigned to whichever set it falls in (or to
\emph{undecided} if none). Let
$n_y$, $n_b$, $n_d$ be the pixel counts for yellow, brown, and dark
respectively. The reported \emph{spot coverage} is
\begin{equation}
\sigma = \frac{n_b + n_d}{\max(n_y + n_b + n_d, \, 1)} \cdot 100\%
\end{equation}
and the reported \emph{stage} is
\begin{equation}
\text{stage}(\sigma) =
\begin{cases}
\text{Fresh}     & \text{if } \sigma <  10\% \\
\text{Spotting}  & \text{if } 10\% \leq \sigma < 40\% \\
\text{Spoiling}  & \text{if } 40\% \leq \sigma < 70\% \\
\text{Spoiled}   & \text{if } \sigma \geq 70\%.
\end{cases}
\end{equation}

The classifier reports both the percentage and the stage label
(Figure~\ref{fig:app}b shows a representative output of
$\sigma = 30.6\%$, labeled \emph{Spotting}). We deliberately avoid
hidden state across frames: each capture is classified independently,
which makes the output reproducible from the captured image alone and
serves as a transparent baseline against which future
machine-learning--based detectors can be compared.

\subsection{Per-Product Risk Score}
\label{sec:algo-risk}

The per-product risk score $R \in [0, 100]$ combines six channels:
expiration time remaining, ambient temperature, ambient humidity, gas
trend, visible spotting (for products supporting visual analysis),
and weight change rate. The score is bounded at $100$ and partitioned
into three bands for display:
\begin{itemize}
  \item \textbf{Fresh} \ for $R < 40$,
  \item \textbf{Consume Soon} \ for $40 \leq R < 70$,
  \item \textbf{Spoilage Risk} \ for $R \geq 70$.
\end{itemize}

Concretely, given a product $p$ with metadata stored in the in-memory
product store, the risk score is computed as
\begin{equation}
R(p) = \min(100, \,
   r_{\text{exp}}(p) + r_T + r_H + r_G + r_V(p) + r_W(p)),
\end{equation}
where each term is a per-channel risk contribution, capped between
$0$ and a per-channel maximum, and rounded down to the nearest integer.
\begin{itemize}
  \item $r_{\text{exp}}$ is $0$ when the expiration date is more than
        a week away, ramps to $40$ when the product reaches its
        expiration date, and pins at $40$ thereafter.
  \item $r_T$ is $\max(0, T - 8^{\circ}\mathrm{C}) \cdot 1.5$,
        capped at $20$.
  \item $r_H$ is $\max(0, H - 70\%) \cdot 0.5$, capped at $10$.
  \item $r_G$ is $\max(0, G - G_{\text{baseline}}) \cdot 0.4$, capped
        at $10$ (where $G_{\text{baseline}}$ is the cabinet's
        recalibration baseline).
  \item $r_V$ is the spot-coverage percentage divided by~5, capped at
        $10$, applied only to products whose category supports visual
        analysis (\texttt{fruit}, in this prototype).
  \item $r_W$ is $5$ if the product's measured weight has dropped by
        more than $50\%$ since registration (indicating that the
        product has been mostly consumed) and $0$ otherwise.
\end{itemize}

Each per-channel maximum is intentionally smaller than $100$ so that
no single sensor channel can dominate the score: a high gas reading
alone, for instance, is not sufficient to flag a product as
\emph{Spoilage Risk}; the risk must be \emph{multi-channel}, which
matches the intuition that food deterioration tends to manifest in
more than one indicator simultaneously.

The score is recomputed inside the Flutter application each time a
sensor reading arrives, which gives the alert engine a current view of
every registered product without persisting state across application
restarts.
